\clearpage
\begin{examplebox}
\textbf{例1｜求负数的取整值。} 求 $\lfloor-3.2\rfloor$ 与 $\{-3.2\}$。因为 $-4\le-3.2<-3$，所以答案为 $-4$，而不是 $-3$。再由小数部分定义，$\{-3.2\}=-3.2-(-4)=0.8$。

\textbf{例2｜解方程。} 解 $\lfloor2x-1\rfloor=3$。
\[3\le2x-1<4\iff4\le2x<5\iff2\le x<\frac52.\]
答案为 $[2,5/2)$。

\textbf{例3｜解不等式。} 解 $\lfloor3-2x\rfloor\le2$。
\[3-2x<3\iff-2x<0\iff x>0.\]
除以负数时不等号反向。
\end{examplebox}
\clearpage
\begin{examplebox}
\textbf{例4｜处理两个取整符号。} 解 $\lfloor x\rfloor+\lfloor-x\rfloor=-1$。由相反数关系，解集为 $\mathbb R\setminus\mathbb Z$。

\textbf{例5｜整数计数。} 求 $1$ 至 $100$ 中能被 $7$ 整除的正整数个数。设倍数为 $7k$，则 $1\le k\le100/7$。故个数为 $\lfloor100/7\rfloor=14$。

\textbf{例6｜非整数边界。} 解 $\lfloor x\rfloor\le2.5$。因左侧是整数，等价于 $\lfloor x\rfloor\le2$，再转化为 $x<3$。

\textbf{例7｜小数部分与整数分解。} 解 $\{x\}=\frac34$。令 $n=\lfloor x\rfloor\in\mathbb Z$，则 $x-n=\frac34$，即 $x=n+\frac34$。反过来，这些数均满足 $n\le x<n+1$，故解集为 $\{n+\frac34\mid n\in\mathbb Z\}$。

\textbf{例8｜小数部分的周期与进位。} 求 $\{2026.7\}$ 与 $\lfloor 2.7+3.6\rfloor$。由整数平移，$\{2026.7\}=\{0.7\}=0.7$。又 $2.7+3.6=5+1.3=6+0.3$，故取整结果为 $6$。

\textbf{变式训练：} （1）$\lfloor 1-3x\rfloor=2$；（2）$\lfloor x\rfloor+\lfloor x+1\rfloor=7$；（3）$\lfloor x\rfloor+\lfloor-x\rfloor=0$。

\textbf{答案：} （1）$(-2/3,-1/3]$；（2）$[3,4)$；（3）$\mathbb Z$。
\end{examplebox}
